\chapter{METODOLOGI}

% Ubah konten-konten berikut sesuai dengan isi dari metodologi

\section{Metode yang digunakan}
Metode yang digunakan pada penelitian ini adalah Vision Transformer yang termasuk dalam \textit{deep learning}.
Dataset yang digunakan adalah citra pola sisik pada kepala penyu yang diperoleh dari berbagai macam sumber fotografi.
Dataset ini kemudian akan dilatih dengan menggunakan Pytorch hingga menghasilkan model yang mampu melakukan Re-Identifikasi.

\subsection{Pengumpulan Dataset}
Pengumpulan dataset dilakukan dengan mencari dataset yang tersedia secara publik di internet. 
Dataset yang dicari adalah citra kepala penyu yang menampilkan pola sisik pada kepalanya.
Dataset dicari sebanyak mungkin dikarenakan sifat Vision Transformer yang membutuhkan banyak data untuk dapat bekerja secara efektif.
Selain itu, dibutuhkan dataset dengan sudut kamera dan pencahayaan yang berbeda-beda agar model yang dihasilkan dapat bekerja pada situasi yang bereagam.

\subsection{Pengolahan Dataset}
Pengolahan dataset dilakukan dengan memproses dataset menjadi data dengan parameter yang diinginkan seperti resolusi yang sesuai.
Pengolahan dataset dilakukan untuk menyiapkan gambar agar dapat digunakan secara optimal dalam pelatihan model.
eknik augmentasi data juga diterapkan untuk meningkatkan keragaman data pelatihan, dengan cara melakukan rotasi, flipping, zooming, dan perubahan warna pada gambar. 
Hal ini bertujuan untuk mencegah overfitting dan meningkatkan kemampuan model dalam mengenali berbagai variasi kondisi gambar.
Dataset akan dibagi menjadi 80\% \textit{training} dan 20\% pengujian.

\subsection{Penyusunan Sistem}
Model yang digunakan dalam penelitian ini adalah Vision Transformer (ViT), yang merupakan arsitektur deep learning terbaru yang dirancang khusus untuk pemrosesan gambar.
Proses ini dilakukan dengan penyusun arsitektur model agar dapat bekerja dengan dataset yang digunakan dan menghasilkan akurasi sesuai yang diharapkan.
Selain itu dapat dilakukan \textit{fine-tuning} model yang telah dilatih sebelumnya dengan dataset besar, seperti ImageNet, sebagai dasar.

\subsection{Pelatihan Model}
Pelatihan model dilakukan dengan menggunakan optimizer AdamW yang terkenal efisien dalam mengoptimalkan model berbasis transformer. 
Learning rate diatur dengan menggunakan teknik scheduler, seperti cosine annealing, untuk menyesuaikan laju pembelajaran selama proses pelatihan.
Fungsi loss yang digunakan adalah Cross-Entropy Loss untuk tugas klasifikasi, atau Triplet Loss untuk pembelajaran embedding, yang memungkinkan model untuk belajar perbedaan antara individu penyu.
Pelatihan model dilakukan menggunakan GPU untuk mempercepat proses komputasi.

\subsection{Evaluasi Model}
Untuk mengevaluasi kinerja model, beberapa metrik digunakan, seperti akurasi klasifikasi yang mengukur sejauh mana model dapat mengenali identitas individu penyu. 
Selain itu, mAP (mean Average Precision) digunakan untuk menilai kemampuan model dalam melakukan re-identifikasi penyu. 
ROC curve dan AUC juga digunakan untuk mengevaluasi kinerja model dalam hal kemampuan deteksi dan klasifikasi. 
Validasi model dilakukan menggunakan dataset yang berbeda untuk mengukur kemampuan generalisasi model terhadap data yang belum pernah dilihat sebelumnya, sehingga dapat dipastikan bahwa model tidak overfit terhadap data pelatihan.

\section{Bahan dan peralatan yang digunakan}
Bahan yang digunakan pada penelitian ini adalah dataset berupa citra kepala penyu yang tersedia secara publik.
Beberapa diantaranya adalah SeaTurtleIDHeads, SeaTurtleID2022, dan SeaTurtleID.
Sedangkan alat yang digunakan pada penelitian ini adalah laptop PC, dan beberapa perangkat lunak pendukung seperti Microsoft Visual Studio Code, Jupyter Notebook, Python, Pytorch, dan Torchvision.
